\chapter{Descrizione dell'Algoritmo}

\section{Idea Generale}

Il problema richiede di trovare, per ogni città $v$, il percorso dalla 
capitale $0$ a $v$ che minimizza la funzione:

\begin{equation}
    f(\pi) = \sum_{e \in \pi} c(e) - \max_{e \in \pi} c(e)
\end{equation}

Questa funzione obiettivo non è decomponibile nel senso classico di 
Dijkstra: il costo effettivo di un percorso dipende dalla tratta massima 
dell'intero percorso, che non è nota fino al completamento del percorso stesso.

La soluzione adottata è una \textbf{variante di Dijkstra} in cui lo stato 
non è semplicemente la distanza dal nodo sorgente, ma una coppia 
$(\texttt{totalCost}, \texttt{maxCost})$ che permette di calcolare il 
costo effettivo in ogni momento.

\section{Variante di Dijkstra}

\subsection{Stato}

Ogni elemento della coda di priorità è uno \texttt{State}:

\begin{lstlisting}
State {
    int totalCost;  // somma dei costi delle tratte percorse
    int maxCost;    // costo della tratta piu' onerosa
    int node;       // nodo corrente
}
\end{lstlisting}

Il costo effettivo di uno stato è \texttt{totalCost - maxCost}.

\subsection{Inizializzazione}

La coda di priorità viene inizializzata con lo stato della capitale:

\begin{lstlisting}
dist[0]      = 0;
totalCost[0] = 0;
maxEdge[0]   = 0;
pq.push({0, 0, 0});
\end{lstlisting}

Tutti gli altri nodi vengono inizializzati con distanza $+\infty$.

\subsection{Rilassamento}

Per ogni arco $(u, v, c)$ estratto dalla coda, il nuovo stato verso $v$ è:

\begin{equation}
    \texttt{newTotal} = \texttt{curr.totalCost} + c
\end{equation}
\begin{equation}
    \texttt{newMax} = \max(\texttt{curr.maxCost},\ c)
\end{equation}
\begin{equation}
    \texttt{newEff} = \texttt{newTotal} - \texttt{newMax}
\end{equation}

Se \texttt{newEff} $<$ \texttt{dist[v]}, lo stato viene aggiornato e 
inserito nella coda:

\begin{lstlisting}
if (newEff < dist[e.dest]) {
    dist[e.dest]      = newEff;
    totalCost[e.dest] = newTotal;
    maxEdge[e.dest]   = newMax;
    parent[e.dest]    = curr.node;
    pq.push({newTotal, newMax, e.dest});
}
\end{lstlisting}

\subsection{Gestione degli Stati Obsoleti}

Poiché la coda di priorità può contenere più stati per lo stesso nodo, 
viene utilizzato il controllo:

\begin{lstlisting}
if (curr.effectiveCost() > dist[curr.node])
    continue;
\end{lstlisting}

Questo scarta gli stati inseriti in precedenza e superati da uno 
migliore, evitando elaborazioni inutili.

\subsection{Terminazione}

L'algoritmo termina quando la coda di priorità è vuota. A quel punto, 
per ogni nodo $v$, \texttt{dist[v]} contiene il costo effettivo minimo 
del percorso ottimo dalla capitale.

\section{Ricostruzione del Percorso}

Il percorso ottimo viene ricostruito tramite il vettore \texttt{parent}, 
che per ogni nodo $v$ memorizza il nodo precedente nel percorso ottimo. 
Il metodo \texttt{reconstructPath} risale il vettore da \texttt{dest} 
fino alla capitale, producendo il percorso in ordine inverso, che viene 
poi ribaltato:

\begin{lstlisting}
vector<int> path;
int v = dest;
while (v != -1) {
    path.push_back(v);
    v = parent[v];
}
reverse(path.begin(), path.end());
\end{lstlisting}

\section{Correttezza}

La correttezza dell'algoritmo si basa sul fatto che la funzione 
\texttt{effectiveCost} è \textbf{monotona non decrescente} lungo 
qualsiasi percorso: aggiungere una tratta non può diminuire il costo 
effettivo, poiché il nuovo \texttt{totalCost} aumenta sempre di almeno 
1 (i costi sono positivi) mentre il nuovo \texttt{maxCost} può solo 
rimanere uguale o aumentare. Questa proprietà garantisce che il primo 
percorso estratto dalla coda per un nodo sia quello ottimo, 
analogamente al Dijkstra classico.

\section{Lettura dell'Input}

Il file di input viene letto in \texttt{main.cpp} tramite 
\texttt{std::ifstream}. La prima riga fornisce $N$ e $P$; le successive 
$P$ righe forniscono le tratte $(S, E, C)$ che vengono aggiunte al 
grafo tramite \texttt{addEdge}, che inserisce ogni tratta in entrambe 
le direzioni essendo il grafo non orientato.